% ============================================================
%  Beamer презентация – Дипломна защита
%  Compile-time ORM библиотека на C++
%  Алекс Цветанов  |  ТУ–София  |  2026
%  Компилация: lualatex presentation.tex (два пъти за ToC)
% ============================================================
\documentclass[aspectratio=169,12pt]{beamer}

% ── Кирилица и езикова поддръжка (LuaLaTeX – native UTF-8) ─
\usepackage{fontspec}
\usepackage{polyglossia}
\setdefaultlanguage{bulgarian}
\setotherlanguage{english}
% Windows system fonts — гарантирано присъстват, пълна Кирилица
\setmainfont{Times New Roman}
\setsansfont{Arial}
\setmonofont[Scale=0.88]{Courier New}
% Polyglossia изисква явна декларация за всяко семейство
\newfontfamily\cyrillicfont{Times New Roman}
\newfontfamily\cyrillicfontsf{Arial}
\newfontfamily\cyrillicfonttt{Courier New}

% ── Тема и цветове (TU-Sofia blue) ─────────────────────────
\usetheme{Madrid}
\definecolor{tusofia}{RGB}{0,86,145}
\definecolor{tusofiadark}{RGB}{0,60,105}
\definecolor{tusofialight}{RGB}{220,235,248}

\setbeamercolor{palette primary}{bg=tusofia,fg=white}
\setbeamercolor{palette secondary}{bg=tusofiadark,fg=white}
\setbeamercolor{palette tertiary}{bg=tusofia!70,fg=white}
\setbeamercolor{palette quaternary}{bg=tusofia,fg=white}
\setbeamercolor{structure}{fg=tusofia}
\setbeamercolor{block title}{bg=tusofia,fg=white}
\setbeamercolor{block body}{bg=tusofialight,fg=black}
\setbeamercolor{alerted text}{fg=red!70!black}

% ── Шрифт за code ──────────────────────────────────────────
\usepackage{listings}
\usepackage{xcolor}

\definecolor{cppblue}{RGB}{0,0,170}
\definecolor{cppgreen}{RGB}{0,120,0}
\definecolor{cppred}{RGB}{170,0,0}
\definecolor{cppgray}{RGB}{110,110,110}

\lstset{
    language=C++,
    basicstyle=\fontsize{6}{7.5}\ttfamily,
    keywordstyle=\color{cppblue}\bfseries,
    stringstyle=\color{cppred},
    commentstyle=\color{cppgreen}\itshape,
    identifierstyle=\color{black},
    showstringspaces=false,
    breaklines=true,
    breakatwhitespace=true,
    columns=flexible,
    keepspaces=true,
    tabsize=4,
    frame=single,
    framerule=0.4pt,
    rulecolor=\color{cppgray},
    xleftmargin=0.4em,
    xrightmargin=0.2em,
    morekeywords={constexpr,requires,concept,auto,nullptr,noexcept,
                  select,insert,update,deleteq,where,order_by,group_by,
                  limit,prepare,execute,field,Placeholder,ph,
                  IndexedPlaceholder,static_assert,using,namespace},
}

% ── TikZ ───────────────────────────────────────────────────
\usepackage{tikz}
\usetikzlibrary{arrows.meta,positioning,shapes.geometric,fit,calc,decorations.pathreplacing}

\tikzset{
    layerbox/.style={draw=tusofia, rounded corners=3pt, fill=tusofialight,
        align=center, text width=3.4cm, font=\footnotesize,
        minimum width=3.8cm, minimum height=0.55cm, inner sep=3pt},
    appbox/.style={layerbox, fill=green!12, draw=green!60!black},
    ormline/.style={-Latex, thick, color=tusofia},
    capbox/.style={draw=orange!70!black, rounded corners=3pt, fill=orange!10,
        align=center, text width=1.7cm, font=\scriptsize,
        minimum width=1.9cm, minimum height=0.55cm, inner sep=3pt},
}

% ── Графики ────────────────────────────────────────────────
\usepackage{graphicx}
\usepackage{booktabs}
\usepackage{tabularx}

% ── Метаданни ──────────────────────────────────────────────
\title[Compile-time ORM за C++]{%
    Compile-time Object-Relational Mapping\\[0.2em]
    библиотека на C++}
\subtitle{Дипломна работа}
\author[Алекс Цветанов]{%
    Алекс Цветанов, фак.~№~121222225\\[0.3em]
    {\footnotesize Научен ръководител: доц.~д-р~инж.~Галя Павлова}}
\institute[ТУ--София]{%
    Технически университет~--~София\\
    Факултет „Компютърни системи и технологии"\\
    Компютърно и софтуерно инженерство}
\date{2026}

% ══════════════════════════════════════════════════════════
\begin{document}
% ══════════════════════════════════════════════════════════

% ── Заглавна страница ──────────────────────────────────────
\begin{frame}
    \titlepage
\end{frame}

% ── Съдържание ─────────────────────────────────────────────
\begin{frame}{Съдържание}
    \tableofcontents
\end{frame}

% ══════════════════════════════════════════════════════════
\section{Мотивация}
% ══════════════════════════════════════════════════════════

\begin{frame}{Проблемът с runtime заявките}
    \begin{columns}[T]
        \begin{column}{0.54\textwidth}
            \textbf{Традиционен подход:}
            \begin{itemize}
                \item Заявките се описват като \alert{низове}
                \item Валидиране едва при изпълнение
                \item Грешни имена на полета → runtime crash
                \item Несъвместими типове → late discovery
                \item Смяна на backend → \alert{пълно пренаписване}
            \end{itemize}
            \vspace{0.3em}
            \textbf{Изследователски въпрос:}\\
            \textit{Може ли компилаторът да поеме ролята на
            първи валидатор на заявките?}
        \end{column}
        \begin{column}{0.43\textwidth}
            \begin{tabularx}{\linewidth}{@{}XX@{}}
                \toprule
                \textbf{Runtime} & \textbf{Compile-time}\\
                \midrule
                Низ заявки &
                    \texttt{constexpr} IR\\
                {Грешка при\newline изпълнение} & {Грешка при\newline \alert{компилация}}\\
                Late discovery &
                    Early detection\\
                \bottomrule
            \end{tabularx}
            \vspace{0.3em}
            \begin{alertblock}{Цел}
                {\small Компилаторът поема ролята
                на първи валидатор
                на заявките.}
            \end{alertblock}
        \end{column}
    \end{columns}
\end{frame}

\begin{frame}{Съществуващи решения и позиция на разработката}
    \begin{columns}[T]
        \begin{column}{0.52\textwidth}
            \begin{tabularx}{\linewidth}{@{}llX@{}}
                \toprule
                \textbf{Решение} & \textbf{Подход} & \textbf{Ограничение}\\
                \midrule
                SOCI      & runtime bind & низ-базирани заявки\\
                ODB       & codegen      & external compiler\\
                sqlpp11   & EDSL         & само SQL\\
                Hiberlite & ORM          & само SQLite\\
                \bottomrule
            \end{tabularx}
        \end{column}
        \begin{column}{0.45\textwidth}
            \begin{block}{Приносът на разработката}
                \begin{itemize}
                    \item[$+$] \texttt{constexpr} query IR без codegen
                    \item[$+$] Multi-backend: SQL \emph{и} NoSQL
                    \item[$+$] Compile-time capability проверки
                    \item[$+$] Формален connector contract
                \end{itemize}
            \end{block}
        \end{column}
    \end{columns}
\end{frame}

% ══════════════════════════════════════════════════════════
\section{Архитектура}
% ══════════════════════════════════════════════════════════

\begin{frame}{Слоеста архитектура на библиотеката}
    \begin{columns}[T]
        \begin{column}{0.42\textwidth}
            \centering
            \begin{tikzpicture}[node distance=0.28cm]
                \node[appbox] (app)
                    {\textbf{Приложен слой}\\ entity + query};
                \node[layerbox, below=of app] (ent)
                    {\textbf{Entity слой}\\
                    \texttt{property}, \texttt{relationship}};
                \node[layerbox, below=of ent] (ir)
                    {\textbf{Query IR слой}\\
                    \texttt{field}, \texttt{Rule}, заявки};
                \node[layerbox, below=of ir] (conn)
                    {\textbf{Connector слой}\\
                    {\scriptsize\texttt{connector\_trait<DB>}}};
                \node[layerbox, fill=gray!10, draw=gray!50,
                    below=of conn] (drv)
                    {\textbf{Driver / runtime}\\
                    SQL, BSON, Cypher, CQL};
                \draw[ormline] (app)  -- (ent);
                \draw[ormline] (ent)  -- (ir);
                \draw[ormline] (ir)   -- (conn);
                \draw[ormline] (conn) -- (drv);
            \end{tikzpicture}
        \end{column}
        \begin{column}{0.55\textwidth}
            \textbf{Ключови архитектурни принципа:}
            \begin{itemize}
                \item Никакъв виртуален интерфейс\\
                    {\footnotesize trait-базиран, не наследствен contract}
                \item Query IR е \texttt{constexpr}\\
                    {\footnotesize значима валидация при компилация}
                \item Capability tags ограничават\\
                    неподдържаните операции
                \item Смяна на backend = смяна само на\\
                    типа на конектора
            \end{itemize}
            \vspace{0.2em}
            \begin{alertblock}{Централен инвариант}
                Приложният код не комуникира директно с драйвера —
                работи само с C++ типове.
            \end{alertblock}
        \end{column}
    \end{columns}
\end{frame}

% ══════════════════════════════════════════════════════════
\section{Реализация}
% ══════════════════════════════════════════════════════════

\begin{frame}[fragile]{Entity модел и Query IR}
    \begin{columns}[T]
        \begin{column}{0.49\textwidth}
            \textbf{Entity описание:}
\begin{lstlisting}
struct User {
  orm::property<int,        "id">    id;
  orm::property<std::string,"name"> name;
  orm::property<double,     "score"> score;
};

template<> struct orm::table_name_trait<User> {
  static constexpr std::string_view
    value = "users";
};
\end{lstlisting}
        \vspace{0.2em}
        {\scriptsize
        \texttt{property<T,"col">} носи тип, compile-time
        name и member pointer — без macro, без codegen.}
        \end{column}
        \begin{column}{0.49\textwidth}
            \textbf{Query IR конструкции:}
\begin{lstlisting}
// field<Ptr> -- wrapper над member ptr
// Rule       -- типизиран предикат
// Placeholder -- анонимен параметър

auto pred =
  orm::field<&User::id>
    == orm::Placeholder<int>{};

// constexpr обект
static constexpr auto q =
  orm::select(
    orm::field<&User::id>,
    orm::field<&User::name>)
  .where(pred);
\end{lstlisting}
        \end{column}
    \end{columns}
\end{frame}

\begin{frame}[fragile,shrink=2]{Изграждане на сложна заявка}
    \begin{columns}[T]
        \begin{column}{0.52\textwidth}
\begin{lstlisting}
using namespace orm::literals;

static constexpr auto q =
  orm::select(
    orm::field<&User::id>,
    orm::field<&User::name>,
    orm::field<&User::score>)
  .where(
    orm::field<&User::score> > 3.14)
  .where(
    orm::field<&User::id>
      == orm::Placeholder<int>{})
  .order_by<orm::order::direction::desc>(
    orm::field<&User::score>)
  .limit(10_per_page & 1_page);

static_assert(
  decltype(q.where_clauses())::size == 2);
\end{lstlisting}
        \end{column}
        \begin{column}{0.45\textwidth}
            \textbf{Характеристики:}
            \begin{itemize}\small
                \item Fluent API — всяка стъпка
                    връща \emph{нов} типизиран обект
                \item \texttt{static constexpr} — IR
                    живее в code сегмента;
                    нулев runtime overhead
                \item \texttt{static\_assert} проверява
                    структурата при компилация
                \item \texttt{group\_by}, \texttt{join<>}
                    — аналогичен синтаксис
            \end{itemize}
            \vspace{0.15em}
            \begin{block}{Pagination UDL}
                {\small\texttt{10\_per\_page \& 2\_page}}\\
                {\scriptsize двата оператора са взаимозаменяеми}
            \end{block}
        \end{column}
    \end{columns}
\end{frame}

\begin{frame}[fragile]{Connector trait и Capability механизъм}
    \begin{columns}[T]
        \begin{column}{0.52\textwidth}
            \textbf{Минимален connector contract:}
\begin{lstlisting}
template<>
struct orm::connector_trait<MyDB> {
  // Задължителни типове
  template<typename T> struct wire_type;
  using cursor_type = ...;

  // Capability тагове
  using supports_joins        = true_type;
  using supports_transactions = true_type;

  // Materialization entry point
  static auto execute(
    MyDB& db, auto query,
    auto&&... params);
};
\end{lstlisting}
        \end{column}
        \begin{column}{0.45\textwidth}
            \textbf{Compile-time capability:}
            \begin{itemize}\small
                \item \texttt{is\_connector<DB>} concept
                    проверява всеки тип
                \item JOIN без \texttt{supports\_joins}
                    → \alert{compile error}
                \item Транзакция без\newline
                    \texttt{supports\_transactions}
                    → \alert{compile error}
            \end{itemize}
            \vspace{0.2em}
            \begin{tikzpicture}[node distance=0.4cm]
                \node[capbox] (ir2) {Query IR};
                \node[capbox, right=0.3cm of ir2] (cap) {Capability?};
                \node[capbox, fill=green!12, draw=green!60,
                    below=0.4cm of cap] (ok) {execute()};
                \node[capbox, fill=red!12, draw=red!60,
                    below=0.4cm of ir2] (err) {compile err};
                \draw[-Latex,thick,color=tusofia] (ir2) -- (cap);
                \draw[-Latex,thick,color=green!60!black] (cap) -- (ok);
                \draw[-Latex,thick,color=red!70] (cap) -- (err);
            \end{tikzpicture}
        \end{column}
    \end{columns}
\end{frame}

% ══════════════════════════════════════════════════════════
\section{Backend сценарии}
% ══════════════════════════════════════════════════════════

\begin{frame}{Multi-backend материализация}
    \begin{columns}[T]
        \begin{column}{0.50\textwidth}
            {\small Един Query IR → различни backend-и:}
            \begin{tabularx}{\linewidth}{@{}lX@{}}
                \toprule
                \textbf{Backend}  & \textbf{Материализация}\\
                \midrule
                SQLite      & \texttt{WHERE id = ?}\\
                PostgreSQL  & \texttt{WHERE id = \$1}\\
                MySQL       & \texttt{WHERE id = ?}\\
                \midrule
                MongoDB     & BSON filter\\
                Redis       & \texttt{GET}/\texttt{SET}/\texttt{DEL}\\
                Neo4j       & Cypher \texttt{MATCH…WHERE}\\
                Cassandra   & CQL \texttt{SELECT…FROM}\\
                \bottomrule
            \end{tabularx}
        \end{column}
        \begin{column}{0.47\textwidth}
            \textbf{Абстракция на връзките:}\\[0.1em]
            \begin{tabularx}{\linewidth}{@{}lX@{}}
                \texttt{embed}     & MongoDB вложен документ\\
                \texttt{reference} & SQL foreign key / JOIN\\
                                   & Redis: key prefix\\
                                   & Neo4j: edge\\
            \end{tabularx}
            \vspace{0.1em}
            \begin{alertblock}{Принцип}
                {\small C++ моделът е единственият
                \emph{логически} носител на схемата.
                Backend-ът определя физическата
                материализация.}
            \end{alertblock}
        \end{column}
    \end{columns}
\end{frame}

\begin{frame}[fragile,shrink=3]{Подготвени заявки}
    \begin{columns}[T]
        \begin{column}{0.52\textwidth}
\begin{lstlisting}
static constexpr auto q =
  orm::select(orm::field<&User::name>)
  .where(orm::field<&User::id>
    == orm::Placeholder<int>{});

auto stmt = db.prepare(q);
auto r1 = stmt.execute(42);
auto r2 = stmt.execute(99);

// Индексирани placeholder-и
using namespace orm;
static constexpr auto q2 =
  orm::select(orm::field<&User::name>)
  .where(orm::field<&User::id>
    == ph<int, _1>)
  .where(orm::field<&User::score>
    > ph<double, _1>);
\end{lstlisting}
        \end{column}
        \begin{column}{0.45\textwidth}
            \textbf{Характеристики:}\\[0.2em]
            {\small
            \textbullet~\texttt{prepare()} →
                \texttt{prepared\_query<Q,DB>}\\[0.2em]
            \textbullet~Многократна употреба без
                повторно рендиране\\[0.2em]
            \textbullet~\texttt{ph<T,\_N>} — индексиран
                placeholder; \texttt{?N} (SQLite)
                или \texttt{\$N} (PostgreSQL)\\[0.2em]
            \textbullet~Един runtime параметър може
                да запълни повече от един слот
            }
            \vspace{0.2em}
            \begin{block}{Верифицирано в}
                {\scriptsize
                \texttt{test\_mockdb.cpp}:\\
                \texttt{PrepareReturnsExecutableObject}\\
                \texttt{IndexedPlaceholderReused\ldots}}
            \end{block}
        \end{column}
    \end{columns}
\end{frame}

% ══════════════════════════════════════════════════════════
\section{Верификация}
% ══════════════════════════════════════════════════════════

\begin{frame}{Верификационна стратегия}
    \begin{columns}[T]
        \begin{column}{0.50\textwidth}
            \begin{block}{Тестова инфраструктура}
                \textbf{214 теста} — 100\% преминаване\\
                {\small Windows/MSVC и Linux/GCC\,13 в Docker}
            \end{block}
            \vspace{0.2em}
            \textbf{Покритие:}
            \begin{itemize}\small
                \item Query IR изграждане и структура
                \item Capability-базирано ограничаване
                \item Подготвени заявки и placeholder-и
                \item Schema миграции (\texttt{migrate<DB>})
                \item Нишкова безопасност
                \item Live integration тестове
            \end{itemize}
        \end{column}
        \begin{column}{0.47\textwidth}
            \textbf{Нива на верификация:}
            \begin{enumerate}\small
                \item \textbf{Compile-time:} \texttt{static\_assert},
                    concept constraints
                \item \textbf{Unit:} изолирани компоненти
                    с MockDB
                \item \textbf{Integration:} SQL rendering,
                    live бази данни
                \item \textbf{Docker CI:} пълен suite на Linux
            \end{enumerate}
            \vspace{0.2em}
            \begin{alertblock}{Ключов принцип}
                {\small Не е достатъчно само „да има тестове“
                или само compile-time checks.
                Необходима е комбинация.}
            \end{alertblock}
        \end{column}
    \end{columns}
\end{frame}

% ══════════════════════════════════════════════════════════
\section{Приноси и заключение}
% ══════════════════════════════════════════════════════════

\begin{frame}{Научни и приложни приноси}
    \begin{enumerate}\setlength{\itemsep}{0.35em}
        \item \textbf{Практически приложим Compile-time ORM модел}\\
            {\small без external codegen, без виртуален интерфейс,
            в стандартен C++23}
        \item \textbf{C++ entity описанието като единствен
            логически носител на схемата}\\
            {\small един модел → таблица / документ / key-value /
            граф / wide-column ред}
        \item \textbf{Формален connector contract за разширяемост}\\
            {\small нов backend без промяна в query API;
            capability tags правят ограниченията видими при компилация}
        \item \textbf{Многостепенна верификационна стратегия}\\
            {\small concept constraints + unit + integration + live}
    \end{enumerate}
\end{frame}

\begin{frame}{Заключение}
    \begin{columns}[T]
        \begin{column}{0.52\textwidth}
            \textbf{Постигнато:}
            \begin{itemize}\small
                \item Query IR като \texttt{constexpr} обект —
                    без низове, без runtime reflection
                \item 7 backend-а: SQLite, PG, MySQL,
                    MongoDB, Redis, Neo4j, Cassandra
                \item 214 теста, Docker CI, два компилатора
                \item Connector contract формализиран
                    и демонстриран с реални конектори
            \end{itemize}
        \end{column}
        \begin{column}{0.45\textwidth}
            \textbf{Бъдещо развитие:}
            \begin{itemize}\small
                \item C++26 static reflection →
                    автом. entity описание
                \item Distributed connection pools
                \item Автом. migration за всички
                    backend-и
                \item Performance benchmarking
            \end{itemize}
            \vspace{0.15em}
            \begin{block}{Основен извод}
                {\small Между runtime ORM и driver-only
                съществува жизнеспособна
                арх. позиция:\\[0.1em]
                \textbf{Compile-time ORM слой в C++.}}
            \end{block}
        \end{column}
    \end{columns}
\end{frame}

\begin{frame}[plain]
    \vfill
    \centering
    {\Large\color{tusofia} Благодаря за вниманието!}\\[0.6em]
    {\large Въпроси?}
    \vfill
\end{frame}

\end{document}
